\documentclass[11pt]{article}
\usepackage[margin=25mm]{geometry}
\usepackage{fontspec}\setmainfont{Latin Modern Roman}
\usepackage{amsmath,amssymb,mathtools}
\usepackage{unicode-math}\setmathfont{Latin Modern Math}
\usepackage{xurl,hyperref}
\hypersetup{colorlinks=true,urlcolor=blue}
\setcounter{secnumdepth}{0}\setlength{\parindent}{0pt}\setlength{\parskip}{5pt}\setlength{\emergencystretch}{3em}
\newcommand{\sref}[2]{\hyperlink{source:#1:#2}{[S#2]}}
\newcommand{\cataloguescope}[1]{\paragraph{Recorded target.}#1}
\begin{document}
\section{Kourovka Notebook Problem 21.111}\label{statement}

\subsection{Definitions and mathematical statement}

\paragraph{Shared notation.}$\mathbb N=\{1,2,\ldots\}$, and $\mathbb Z,\mathbb Q,\mathbb R,\mathbb C$ denote the integers, rationals, reals and complex numbers. Any other conventions are specified locally.


For a finite nonabelian simple group $S$, put $G_x=\langle x,\operatorname{Inn}(S)\rangle$ for $1\ne x\in\operatorname{Aut}(S)$. Such an $x$ is a $\tau$-automorphism when
\[
 3\nmid\bigl|\langle g^{-1}xg,h^{-1}xh\rangle\bigr|
 \qquad\text{for every }g,h\in G_x.
\]
A $\tau$-group is an $S$ admitting at least one such automorphism. Exclude the Suzuki simple groups ${}^2B_2(q)$, where $q=2^{2m+1}$ with $m\geq1$. The order of an automorphism is its order as an element of $\operatorname{Aut}(S)$.
\paragraph{Statement recorded in the supplied source.}
Among finite nonabelian simple groups not isomorphic to any ${}^2B_2(q)$:

(a) classify all $\tau$-groups up to isomorphism;

(b) determine whether a $\tau$-automorphism of odd order exists. The pairwise condition concerns every pair of conjugates in $G_x$, not just a selected generating pair. \sref{statement}{2}
\subsection{Short English statement}
Classify the finite simple groups, apart from Suzuki groups, admitting the source’s pairwise 3-divisibility-avoiding automorphisms, and ask whether any such automorphism has odd order.
\subsection{Sources}
\begin{description}
\item[\hypertarget{source:statement:1}{S1}] ULAM.AI and the UnsolvedMath Contributors, \emph{UnsolvedMath: A Curated Collection of Open Mathematics Problems}, catalogue entry. CC BY 4.0. \url{https://huggingface.co/datasets/ulamai/UnsolvedMath}.
\item[\hypertarget{source:statement:2}{S2}] D. O. Revin and N. Y. Yang, Problem 21.111, in E. I. Khukhro and V. D. Mazurov (eds.), \emph{The Kourovka Notebook: Unsolved Problems in Group Theory}, 21st edition, arXiv:1401.0300v48 (6 October 2026). \url{https://arxiv.org/abs/1401.0300v48}.
\end{description}
\label{end:statement}
\end{document}
